%*****************************************
\chapter{Lillama: Large Language Models Compression via Low-Rank Feature Distillation}\label{ch:lillama}
%*****************************************

\section{Introduction}

\begin{figure}[ht]
        \centering
        \includegraphics[width=\linewidth]{papers/lillama/overview.pdf}
        \caption{\textbf{Lillama approach}: \textbf{STEP 1} selects layers to compress for a target compression ratio (e.g., N\%) using various strategies (see Section \ref{sec:algo}). \textbf{STEP 2} compresses and initializes the chosen parameters via SVD. \textbf{STEP 3} distills the low-rank weights with a small calibration dataset.}
        \label{fig:overview}
\end{figure}


\begin{table*}[h!]
\setlength{\tabcolsep}{2pt}
\begin{tabular}{>{\raggedright\arraybackslash}p{0.32\textwidth}|>{\centering\arraybackslash}p{0.12\textwidth}|>{\centering\arraybackslash}p{0.12\textwidth}|>{\centering\arraybackslash}p{0.12\textwidth}|>{\centering\arraybackslash}p{0.12\textwidth}|>{\centering\arraybackslash}p{0.12\textwidth}}
\toprule
\textbf{Method} & \textbf{No Custom Kernel} & \textbf{No Ampere-only} & \textbf{Memory Gain} & \textbf{Is One-Shot} & \textbf{Calib. Data} \\  \hline
SparseGPT \citep{frantar2023sparsegptmassivelanguagemodels} & \xmark & \xmark & \xmark & \cmark & - \\
Wanda \citep{wanda} & \xmark & \xmark & \xmark & \cmark & - \\
Teal \citep{TEAL} & \xmark & \xmark & \xmark & \cmark & - \\
ShearedLLaMA \citep{xia2024sheared} & \cmark & \cmark & \cmark & \xmark & 52B\\ 
Minitron \citep{minitron} & \cmark & \cmark & \cmark & \xmark & 94B \\
\rowcolor[HTML]{EEE0CC}\textbf{Lillama} & \cmark & \cmark & \cmark & \cmark & \textbf{0.013B} \\
\bottomrule
\end{tabular}
\caption{Positioning of our approach within the compression literature. \textbf{No Custom Kernel} indicates that a custom GPU kernel is not required to observe speedup and \textbf{No Ampere-only} indicates that speedup is not restricted to Ampere GPUs. \textbf{Calib. Data} is the total number of tokens in billions used for compression.}
\label{tab:litt}
\end{table*}


As discussed earlier, compressing large language models is a crucial challenge that can be addressed through several complementary approaches: quantization, pruning, and distillation. In this section, we introduce a new one-shot compression method that locally distills low-rank weights and satisfies the following three objectives:
\begin{enumerate}
\item The compression algorithm must be compute-efficient, e.g., it should require minimal computational resources and run in a reasonable amount of time. For instance, iterative magnitude pruning as used in prior work~\citep{frankle2019lottery} is too costly to be applied to large-scale LLMs.
\item It must achieve fast convergence using as little training data as possible; this objective is difficult to meet with many model pruning methods that require significant retraining after pruning to recover the lost performance.
\item The method should not cause a significant drop in performance; ideally, all the capabilities of the original LLM should remain intact.
\end{enumerate}
To address the first constraint, we compress the model with a local distillation objective, as opposed to a global objective that is more costly. We further reduce the cost by compressing and distilling only a subset of all layers.
To address the second constraint, we initialize the compressed layers with Singular Value Decomposition (SVD), which enables
reducing the required number of gradient steps as well as the calibration dataset size. We further improve convergence by combining Teacher and Student activations with a joint distillation loss.
We experimentally show the robustness of our approach by compressing several state-of-the-art small and large Transformers, Mixture-of-Experts, and Mamba LLMs.
Figure~\ref{fig:overview} gives an intuitive overview of the proposed method, which is detailed in Section~\ref{sec:algo} and validated in Section~\ref{sec:xp}.

Our approach addresses some limitations of the compression approaches
presented in Section \ref{sec:llm_compression}. Compared to unstructured pruning
and semi-structured pruning, which result in models with sparse weights and do not achieve speedup without specially designed hardware-specific kernels, our approach, which shrinks
the LLM weights, saves memory across all hardware and leads to speedup, as fewer computations are performed. Structured pruning approaches can be costly: for example, the pruning step of Sheared-LLaMA \citep{xia_sheared_2023} is five times as expensive as standard LM training, according to the authors. In contrast, thanks to local gradient updates, our approach is computationally efficient, allowing us to compress a 47B model within minutes on a single A100 GPU.

There are also some low-rank decomposition approaches that compress models by approximating the pretrained matrices with lower-dimensional ones. Building on the observation that neural networks have low intrinsic dimensions~\citep{li2018measuring}, subsequent work shows that Transformer language models also require lower intrinsic dimensions depending on the language task~\citep{aghajanyan2020intrinsic}. However, other studies show that the weights in some pretrained Transformer models are full-rank, unlike their activations, and propose activation-aware compression methods~\citep{yu_compressing_2023,chen_drone}. Our work is related to these approaches, as we also decompose the matrices of the models. However, we apply a lightweight layer-wise feature distillation objective to better recover from compression and propose different distillation strategies to accelerate convergence and improve performance.

Standard knowledge distillation is data-intensive when the student model's parameters are randomly initialized. Hence, this process is as costly as a standard pretraining stage. Layer-wise distillation can be more suitable in some cases, as intermediate layers may contain certain information, such as speaker identity or phonemes in pretrained speech models \citep{baevski2020wav2vec20frameworkselfsupervised}, \citep{pasad2022layerwiseanalysisselfsupervisedspeech}, \citep{chang2022distilhubertspeechrepresentationlearning}, \citep{chang2021explorationselfsupervisedpretrainedrepresentations}. We leverage the observation that layer activations are low-rank and approximate them with fewer parameters by distilling the Transformer layer activations into low-rank weights
initialized from the SVD of the base weights. This layer-wise distillation approach
enables us to perform local gradient updates, hence making the distillation process cheaper.

\section{Motivation: Low-Rank Activations}
\label{ssec:srank}
A motivation for compressing model weights with SVD would be a low-rank structure in those weights. However, previous studies \citep{chen_drone}, \citep{yu_compressing_2023} have shown that the weights of pretrained Transformer models are often full-rank, while their activations have lower rank.
We empirically confirm this observation on three recent models: Falcon-Mamba, Mistral-v0.1 7B, and Phi-2 3B. There are different measures of rank. In this study, we use a simpler one: stable rank ($srank$) as a proxy measure for rank. The stable rank of a matrix is defined as the following ratio:
\begin{equation}
\text{srank}(X) = \frac{\sum_{i=1}^{r} \sigma_i^2}{\sigma_1^2}
\end{equation}
with $X \in \mathbb{R}^{d_{1} \times d_{2}}$ being the input matrix (activations or weights), $r = \min(d_{1}, d_{2})$, and $\sigma_1 \geq \dots \geq \sigma_r$ the singular values of $X$.

Figure~\ref{fig:srank} shows the stable rank of various weight matrices and layer activations for three models: Falcon-Mamba, Mistral-v0.1 7B, and Phi-2 3B. For all models, activations are significantly lower-rank than weight matrices, which confirms previously published results~\citep{yu_compressing_2023}.

This not only motivates the use of activations for compression, but also explains why applying SVD to weights alone does not work well, as we will show in the experiments. How can we explain the fact that weights are full-rank while activations are low-rank? This is an open question. We know that modern LLMs are pre-trained on large and diverse datasets from the web that contain a great deal of knowledge. Hence, the weights of the model need to be full-rank to capture all this knowledge. However, for a given input from a specific task, only a subset of this knowledge is needed, which may explain the low-rank structure of activations.

\begin{figure*}[h!]
    \centering
    \includegraphics[width=\textwidth]{papers/lillama/stablerank_40p.pdf}
    \caption{\textbf{Activations are low-rank.} Comparison of the stable rank of weights (\texttt{*\_proj}) and layer activations (\texttt{layer\_activations}). Each bar is the average stable rank across all layers.}
    \label{fig:srank}
    \begin{insightbox}{Observation 1: Transformer layer activations are extremely low-rank}
    Transformer layer activations have extremely low stable rank compared to the weights. This extends the previous observations discussed above, which focused on linear activations, to the outputs of entire Transformer layers. Interestingly, the key projection weights also appear to have low stable rank.
    \end{insightbox}
\end{figure*}

\section{Background: Low-Rank Approximation of Transformer Models}
Before presenting our approach, we first introduce two methods for low-rank approximation of Transformer models: low-rank approximation from weights and from features.
\subsection{Low-Rank Approximation from Weights}
Given $W \in \mathbb{R}^{d_{1} \times d_{2}}$, a pretrained weight matrix, and an example $x \in \mathbb{R}^{d_{2}}$, the activations $y \in \mathbb{R}^{d_{1}}$ are computed as $y = Wx$. It is possible to reduce the computations in this operation by using a low-rank approximation of $W$: $y \simeq \Delta Wx = ABx$ with $\Delta W = AB$ the low-rank decomposition of $W$, composed of $A \in \mathbb{R}^{d_{1} \times r}$ and $B \in \mathbb{R}^{r \times d_{2}}$. When the rank $r$ is small enough, the number of parameters $r(d_1+d_2)$ in the low-rank equation is smaller than $d_1d_2$ in the full-rank equation. Estimating the low-rank matrix $\Delta W = AB$ can be formulated as a minimization problem to find the low-rank $\Delta W$ that best approximates $W$:

\begin{equation}
\widehat{\Delta W} = \underset{{\Delta W}}{\mathrm{argmin}} \;\; \|W - {\Delta W}\|_{F}
\end{equation}

It is well-known that this minimization problem can be approached using SVD, a low-rank approximation method that offers the optimal rank-$r$ approximation of a given matrix with respect to the Frobenius norm $\|\cdot\|_{F}$.
SVD decomposes a matrix $W$ into three matrices: $W = USV^{T}$, where $S\in \mathbb{R}^{d_{1} \times d_{2}}$ is a diagonal matrix containing the singular values of $W$ sorted in descending order, and $U \in \mathbb{R}^{d_{1} \times d_{1}}$ and $V \in \mathbb{R}^{d_{2} \times d_{2}}$ are orthonormal matrices. The optimal low-rank approximation of $W \in \mathbb{R}^{d_{1} \times d_{2}}$ can be obtained by keeping the first $r$ singular values, with $r < \min(d_{1}, d_{2})$:

\begin{equation}
\widehat{\Delta W} = (U_{:, :r}S_{:r, :r})V_{:r,:} = AB
\label{eq:svd}
\end{equation}

with $A=U_{:,:r}S_{:r,:r}$ and $B=V_{:r,:}$ where the notation $_{:a, :b}$ refers to the slicing operation.

\subsection{Low-Rank Approximation from Features}
\label{ssec:pca}

As discussed previously, many studies \citep{chen_drone}, \citep{yu_compressing_2023} have shown that the activations (i.e., features) of pretrained transformers have lower rank than the weights. Recent work shows that Transformer activations can be sparsified up to 60\%~\citep{TEAL}.
This motivates using the activations of the linear layers to estimate the low-rank matrices.

Let $\mathcal{D} = \{x_i \in \mathbb{R}^{d_{2}}\}_{1\leq i\leq N}$ be a calibration dataset, and $W \in \mathbb{R}^{d_{1} \times d_{2}}$ the parameters of a linear layer. Finding the low-rank matrix $\Delta W$ that best reproduces the activations $\{y=Wx\}~~~ \forall x\in \mathcal{D}$ involves solving:
\begin{equation}
\widehat{\Delta W} = \underset{{\Delta W}}{\mathrm{argmin}} \;\; \frac{1}{N}\sum\limits_{x \in \mathcal{D}}\|Wx - {\Delta Wx}\|_{F}
\end{equation}
An analytic solution to this minimization problem is the eigendecomposition of the covariance matrix of the activations. We first collect all activations $\mathcal{Y}=\{y=Wx\}_{\forall x\in \mathcal{D}}$. The covariance matrix of the activations can be estimated as $\Sigma = \underset{y \in \mathcal{Y}}{\mathbb{E}}\left[yy^T\right] - \mathbb{E}[y]\mathbb{E}[y]^T$ with $\Sigma \in \mathbb{R}^{d_{1} \times d_{1}}$. Since $\Sigma$ is a diagonalizable matrix, we can apply its eigendecomposition: $\Sigma = USU^{T}$, where $U \in \mathbb{R}^{d_{1} \times d_1}$ contains the eigenvectors of $\Sigma$ and $S \in \mathbb{R}^{d_{1} \times d_{1}}$ is the diagonal matrix containing the eigenvalues sorted in decreasing order. As for the SVD equation above, we can keep the eigenvectors corresponding to the $r$ largest eigenvalues, which gives us $A = U_{:, :r}$ and $B = U_{:, :r}^{T}W$.

Compared to the low-rank matrices in the SVD equation above, the weights $A$ and $B$ here are learned to reproduce the activations of the base weight $W$, thanks to the minimization objective above. In the next section, we highlight the limitations of this approach and introduce our proposed method.

\section{Proposed Approach}
\begin{figure}[t]
    \centering
    \subfloat[Teacher\label{fig:distillation_teacher}]{%
        \includegraphics[width=0.30\linewidth]{papers/lillama/Teacher.pdf}}\hfill
    \subfloat[Student\label{fig:distillation_student}]{%
        \includegraphics[width=0.30\linewidth]{papers/lillama/Student.pdf}}\hfill
    \subfloat[Teacher+Student\label{fig:distillation_joint}]{%
        \includegraphics[width=0.30\linewidth]{papers/lillama/Teacher+Student.pdf}}
    \caption[Accelerating convergence through a joint loss]{\textbf{Accelerating convergence by learning from \protect\colorbox{insightframe}{Teacher} and \protect\colorbox{insightback}{Student} activations through a joint loss.} We compare three distillation strategies: \textbf{(a) Teacher}: the input to the compressed student layer comes from the output of the previous teacher layer; \textbf{(b) Student}: the input to the compressed student layer comes from the output of the previous student layer; \textbf{(c) Teacher+Student}: the compressed student layer is applied separately to the output of the previous teacher layer and to the output of the previous student layer, and the joint loss is the sum of the two corresponding distillation losses.}
    \label{fig:distillaton}
\end{figure}
We next identify and address three potential limitations of the approach presented in Section \ref{ssec:pca}.

\subsection{Non-linear Feature Approximation.}
The analytical solution above is limited to activations from linear layers.
To generalize to non-linear modules, we first observe that this objective can be seen as a feature distillation objective that may be optimized numerically by gradient descent rather than analytically by eigendecomposition. Given the input batch $X \in \mathbb{R}^{d \times b}$ of $b$ examples, we denote the output activations of the $i^{th}$ Teacher module $\mathcal{T}^{(i)}$ as
\begin{equation}
Y^{(i)}= \mathcal{T}^{(i)}(X; \Theta^{(i)})
\end{equation}
where $\Theta^{(i)}$ are the original pretrained matrices of the $i^{th}$ Teacher, and $Y^{(i)} \in \mathbb{R}^{d \times b}$ are its output activations. Similarly, we denote the output activations of the $i^{th}$ Student module $\mathcal{S}^{(i)}$ as:
\begin{equation}
\widehat{Y}^{(i)}= \mathcal{S}^{(i)}(X; \Delta \Theta^{(i)})
\end{equation}
where the Student module is parameterized with the low-rank matrices $\Delta \Theta^{(i)}$ and $\widehat{Y}^{(i)} \in \mathbb{R}^{d \times b}$ are the output activations. We can express the distillation objective of the $i^{th}$ Student module as:
\begin{equation}
\widehat{\Delta  \Theta^{(i)}} = \underset{{\Delta  \Theta^{(i)}}}{\mathrm{argmin}} \;\; \mathcal{L}^{(i)}(Y^{(i)}, \; \widehat{Y}^{(i)})
\end{equation}
where $\widehat{\Delta  \Theta^{(i)}}$ are the estimated low-rank matrices of the $i^{th}$ Student module and $\mathcal{L}^{(i)}(Y^{(i)}, \; \widehat{Y}^{(i)})$ is the loss measuring the distance between the activations of the Teacher and the Student.
We use the combined loss from prior work~\citep{chang2022distilhubertspeechrepresentationlearning} because we observed instability with the $\ell_1$ loss alone:
\begin{align}
\mathcal{L}^{(i)} &= \mathcal{L}^{(i)}_{\ell_1} + \mathcal{L}^{(i)}_{\cos} \\
&= \sum_{t=1}^{b} \left[ \frac{1}{D} \left\| Y^{(i)}_{t} - \widehat{Y}^{(i)}_{t} \right\|_1
- \log \sigma \left( \cos \left( Y^{(i)}_{t}, \widehat{Y}^{(i)}_{t} \right) \right) \right] \notag
\end{align}
where $D$ is the hidden-vector dimension, $\sigma$ is the sigmoid activation and cos$(\cdot, \cdot)$ is the cosine similarity. We propose to approximate the low-rank weights of the Students through gradient descent. In the linear case, this should converge towards the eigendecomposition solution described in Section \ref{ssec:pca}.
However, as we show later, we extend the distillation process to non-linear modules, which requires numerical optimization. We also show that initializing the Student's low-rank parameters with SVD, rather than randomly, improves convergence.

\subsection{Beyond Teacher-only Activations.}
The second potential limitation of the approach described in Section \ref{ssec:pca} and equations above is that the Student and the Teacher modules take the same input $X$, which is the output of the previous layer of the Teacher. This means that the Student is trained from the activations of the Teacher module only, which are not available at inference time. This approach has been used to approximate the linear layers of Transformer models under the name Atomic Feature Mimicking~\citep{yu_compressing_2023}. Let $\widehat{Y}^{(i)}_{T} = \mathcal{S}^{(i)}(Y^{(i-1)}_{T}; \Delta \Theta^{(i)})$ be the output activations of the Student module when fed $Y^{(i-1)}_{T}$, the output activations of the Teacher at the previous layer $i-1$. The loss can then be written as:
\begin{equation}
\mathcal{L}^{(i)}_{T} = \mathcal{L}^{(i)}(Y^{(i)}_{T}, \; \widehat{Y}^{(i)}_{T})
\end{equation}
where $Y^{(i)}_{T}$ are the gold activations of the current Teacher. This approach can lead to fast convergence as the Student modules learn from gold and high-quality activations produced by the Teacher modules. However, it can cause a performance drop during inference as the Teacher activations are not available. An alternative approach is to have the Student module take as input the activations of previous Student modules rather than those of Teacher modules:
\begin{equation}
\mathcal{L}^{(i)}_{S} = \mathcal{L}^{(i)}(Y^{(i)}_{T}, \; \widehat{Y}^{(i)}_{S})
\end{equation}
where $\widehat{Y}^{(i)}_{S} = \mathcal{S}^{(i)}(\widehat{Y}^{(i-1)}_{S}; \Delta \Theta^{(i)})$ denotes the output activations of the current Student module when taking as input the output activations $\widehat{Y}^{(i-1)}_{S}$ of the Student module at the previous layer, and $Y^{(i)}_{T}$ are the gold activations of the corresponding Teacher module.
However, although we found that this loss generally performs
better than the previous one, it leads to slow convergence due to
the errors in the activations of the Student modules. To address both issues,
we propose to combine the two losses:
\begin{equation}
\mathcal{L}^{(i)}_{T+S} = \mathcal{L}^{(i)}_{T} + \mathcal{L}^{(i)}_{S}
\end{equation}

This distillation approach is illustrated in Figure~\ref{fig:distillation_joint}. We don't introduce any hyperparameter for controlling the importance of one loss over another in the joint loss. We leave this for future work.

\subsection{Module-wise Distillation.}
Finally, we observe that the modules need not be limited to linear layers: any sub-part of the model that outputs low-rank activations can be compressed with our method. In our experiments, we distill at the Transformer layer level and leave comparisons with other distillation levels for future work. We also experiment with models other than pretrained Transformer-based language models.

\section{Fast and Memory-Efficient Compression}
\label{sec:algo}

Computing the optimal compression rank of each matrix individually may be costly: while search algorithms can be a solution for smaller models, they can be difficult to scale to LLMs with billions of parameters.
Assuming a target compression rate $N$, we consider three simple strategies: \textbf{uniform}, \textbf{top}, and \textbf{bottom}-first compression. Figure~\ref{fig:strategies} summarizes the strengths and weaknesses of each strategy.

\begin{figure}[h!]
    \centering
    \includegraphics[width=\linewidth]{papers/lillama/strategies.pdf}
    \caption{Illustration of 20\% compression when using the three strategies: bottom, top, and uniform. Green boxes are compressed layers with low-rank matrices while gray boxes are non-compressed ones. $l_{1}$, $l_{2}$, ..., $l_{n-1}$, $l_{n}$ are the layer indices.}
    \label{fig:strategies}
\end{figure}

\paragraph{Uniform Weights Compression}
This approach removes $N$\% of the parameters from each weight matrix in the model. It makes no assumptions about the importance of different layers but is costly because it requires computing the low-rank approximation of all weight matrices. It also requires training all the model's layers, increasing GPU memory and computation requirements through forward and backward passes across the entire model.

\paragraph{Top Layers First Compression}
This approach removes $N$\% of the parameters from the model by prioritizing the top layers. This is more efficient than the uniform approach as it requires computing the low-rank approximation of fewer weight matrices and training fewer layers. However, it still requires loading the full model in GPU memory and performing a forward pass through all the layers, although only a subset is compressed and trained.

\paragraph{Bottom Layers First Compression}
This approach removes $N$\% of the parameters from the model by prioritizing the bottom layers. Compared to the previous two approaches, this one is more efficient as it requires computing the low-rank approximation of fewer weight matrices and training fewer layers. It also allows us to load only a subset of the model into GPU memory and run forward passes through only the first compressed layers, which is not the case for the previous two approaches. As we will show, this approach can compress large models that could not fit into GPU memory. Algorithm~\ref{alg:two} describes the bottom-first compression approach.

\begin{algorithm}[!ht]
\small
\caption{Bottom Layers First Compression Algorithm}\label{alg:two}
\KwIn{$\mathcal{M}$ is the base model; $S$ is the target size of the compressed model; $k$ is the minimum possible rank to set; $m$ is the increment when generating ranks}
\KwResult{Low-rank compressed $\mathcal{M}$}
$\mathcal{R} \gets \emptyset$\;
\For{each weight matrix $W \in \mathbb{R}^{d_{1} \times d_{2}}$ in $\mathcal{M}$}{
    $b \gets \min(d_{1}, d_{2})$\;
    \For{$r = k;\ r \leq b;\ r = r + m$}{
        \If{$r \times (d_{1} + d_{2}) < d_{1} \times d_{2}$}{
            $\mathcal{R} \gets \mathcal{R} \cup \{(r, W)\}$\;
        }
    }
}
sort $\mathcal{R}$ by layer index in increasing order, then by rank in decreasing order\;
\While{$|\mathcal{M}| > S$}{
    get the next $(r, W)$ from $\mathcal{R}$\;
    compute $A, B$ from $W$ using Eq.~\ref{eq:svd} with rank $r$\;
    replace $W$ by $A, B$ in $\mathcal{M}$\;
}
\Return{$\mathcal{M}$}\;
\end{algorithm}

\subsection{Ease of Implementation}
The approach is straightforward to implement. Figure~\ref{fig:pseudocode_impl} shows a simple pseudocode implementation that can be easily integrated into most recent PyTorch models.

\begin{figure}[h!]
    \centering
    \includegraphics[width=\linewidth]{papers/lillama/pseudo-code.pdf}
    \caption{Example of a PyTorch implementation of our approach using the Teacher-only distillation strategy.}
    \label{fig:pseudocode_impl}
\end{figure}

\section{Experiments on Transformers}
\label{sec:xp}
\subsection{Setup}
\paragraph{Models.}
We evaluate our method on several Transformer-based LLMs: a 47B Mixture-of-Experts model (Mixtral-8x7B-v0.1), medium-sized LLMs (Mistral-v0.1 7B and Phi-3 14B), and a smaller model (Phi-2 3B).

\paragraph{Data.}
For calibration, we use 13 million tokens randomly sampled from Slim-Orca \citep{SlimOrca}, an open-source replication of Orca \citep{mukherjee2023orca}. In preliminary experiments, we also tested other datasets such as RedPajama \citep{together2023redpajama}, but observed better results with instruction data.

\paragraph{Implementation Details.}
We implement our approach using version 4.44.2 of the Hugging Face Transformers library~\citep{wolf2020huggingfaces} and PyTorch version 2.4.0+cu121. For all experiments, we use the AdamW optimizer with a learning rate of $8.6e^{-4}$. Note that since our approach is local, each student layer has its own optimizer. This also allows us to perform local gradient updates and avoid storing the entire large PyTorch computation graph in memory. All experiments were conducted on a single A100 GPU, regardless of the model.
When compressing the models using Algorithm~\ref{alg:two}, we set a minimum rank of $k=1024$ for all models except Phi-3 14B, for which we found higher ranks work well, so we set $k=1536$. For all models, we set $m=256$ for the increment to generate ranks.
We apply the Teacher+Student loss, as illustrated in Figure~\ref{fig:distillation_joint}, to all models, as we found that this loss converges faster and generally performs better. Additionally, we use the bottom-first compression strategy, which proved to be effective and less memory-intensive than other strategies.

\paragraph{Evaluation.}
To measure inference speed, we evaluate \textit{prefill throughput} by timing a batch of prompts passed through the entire model and reporting the results in tokens per second. Using lm-evaluation-harness \citep{eval-harness}, we evaluate zero-shot performance on 9 downstream tasks: TruthfulQA (TruQA) \citep{lin-etal-2022-truthfulqa}, Social IQA (SIQA) \citep{sap2019social}, LogiQA \citep{liu2020logiqa}, WinoGrande (WinoG) \citep{sakaguchi2019winogrande}, ARC Easy (ARC-E) \citep{arc}, ARC Challenge (ARC-C) \citep{arc}, BoolQ \citep{boolq}, PIQA \citep{piqa}, and OpenBookQA (OBQA) \citep{obqa}. We compare non-compressed (0\%) and compressed variants.
We provide in Table~\ref{tab:task_metric_map} the metrics we used to evaluate the models in Section \ref{sec:xp}.
\begin{table}[!ht]
\centering
    \begin{tabular}{ll}
    \toprule
    \textbf{Task} & \textbf{Metric} \\ 
    \midrule
    arc\_challenge   & acc\_norm \\ 
    arc\_easy        & acc\_norm \\ 
    piqa             & acc\_norm \\ 
    social\_iqa      & acc \\ 
    logiqa           & acc\_norm \\ 
    truthfulqa\_mc2  & acc \\ 
    winogrande       & acc \\ 
    boolq            & acc \\ 
    openbookqa       & acc\_norm \\ 
    \bottomrule
    \end{tabular}
\caption{Benchmarks and corresponding metrics that were used to evaluate the models.}
\label{tab:task_metric_map}
\end{table}

\subsection{Results}
\paragraph{Zero-shot evaluation under different compression ratios.}
Table~\ref{tab:main_result} reports zero-shot results for base models and 20\%, 25\%, and 30\% compressed models, showing that the method remains robust as compression increases.

\paragraph{Low-rank models are lighter and faster.}
Table~\ref{tab:mem-speed} shows memory savings (VRAM) and speedup of compressed models. We can see that 20\% compressed models are lighter and up to 20\% faster.

\paragraph{Our approach can efficiently create small language models.}
We compressed Phi-2 3B to 1.7B (40\% parameter reduction) and compare its performance to recent models of similar size in Table~\ref{tab:smollm}.

\paragraph{Compressed models recover well with fine-tuning.}
Table~\ref{tab:ft_result} shows that fine-tuning these compressed models leads to performance recovery.

\begin{insightbox}{Observation 2: Low-rank compressed LLMs can continue to improve}
As shown by the Mistral fine-tuning results in Table~\ref{tab:ft_result}, the bottleneck representation created by the low-rank structure does not seem to prevent the model from improving. With longer fine-tuning, it could potentially recover most of the lost performance.
\end{insightbox}

\begin{table*}[t]
    \resizebox{\textwidth}{!}{%
    \setlength{\tabcolsep}{3pt}
    \begin{tabular}{cccccccccccc}
        \toprule
        \textbf{Model} & \textbf{Reduction} & \textbf{TruQA} & \textbf{SIQA} & \textbf{LogiQA} & \textbf{WinoG} & \textbf{ARC-E} & \textbf{ARC-C} & \textbf{BoolQ} & \textbf{PIQA} & \textbf{OBQA} & \textbf{Average} \\
        \midrule 
        \multirow{4}*{\centering Phi-3 14B}
        & 0\% & 57.63 & 57.06 & 37.94 & 76.01 & 81.36 & 61.60 & 88.56 & 81.34 & 50.40 & \textbf{65.77} \\
        & 20\% & 57.88 & 50.26 & 34.10 & 72.53 & 83.54 & 59.22 & 85.32 & 80.30 & 49.20 & \textbf{63.59} \\
        & 25\% & 56.37 & 51.38 & 32.26 & 70.56 & 80.98 & 57.51 & 83.88 & 79.76 & 47.00 & \textbf{62.19} \\
        & 30\% & 55.32 & 50.41 & 31.18 & 70.32 & 78.70 & 55.55 & 84.10 & 76.99 & 45.00 & \textbf{60.98} \\
        \midrule
        \multirow{4}*{\centering Phi-2 3B}
        & 0\% & 44.40 & 55.42 & 30.57 & 76.16 & 78.16 & 54.18 & 83.21 & 79.11 & 51.20 & \textbf{61.38} \\
        & 20\% & 46.04 & 52.56 & 29.80 & 71.51 & 72.81 & 46.50 & 75.05 & 76.55 & 45.60 & \textbf{57.38} \\
        & 25\% & 44.55 & 51.64 & 30.88 & 71.35 & 71.93 & 45.99 & 79.79 & 74.43 & 43.20 & \textbf{57.08} \\
        & 30\% & 44.12 & 51.59 & 30.41 & 70.72 & 69.99 & 44.37 & 78.17 & 74.54 & 43.20 & \textbf{56.35} \\
        \midrule
        \multirow{4}*{\centering Mistral-v0.1 7B}
        & 0\% & 42.60 & 46.57 & 29.80 & 73.80 & 79.55 & 53.92 & 83.70 & 82.10 & 44.00 & \textbf{59.56} \\
        & 20\% & 44.14 & 45.50 & 28.26 & 66.69 & 73.11 & 46.33 & 77.00 & 77.37 & 41.20 & \textbf{55.51} \\
        & 25\% & 40.52 & 44.93 & 30.11 & 68.11 & 65.07 & 40.61 & 78.44 & 73.99 & 40.20 & \textbf{53.55} \\
        & 30\% & 40.08 & 44.37 & 30.57 & 66.93 & 65.53 & 40.61 & 77.98 & 73.18 & 39.00 & \textbf{53.14} \\
        \bottomrule
    \end{tabular}}
    \caption{\textbf{Zero-shot evaluation results} for the base non-compressed models (0\%), and the 20\%, 25\%, and 30\% compressed models.}
    \label{tab:main_result}
\end{table*}

\begin{table*}
    \resizebox{1.06\linewidth}{!}{%
    \begin{tabular}{cccc|ccccccc}
    \toprule
       \textbf{Model} & \textbf{Reduction}  & \textbf{Model Size} & \textbf{VRAM} &  s = 512 & s = 1024 &  s = 2048 & s = 4096 & s = 8192 & s = 16384 \\
    \midrule
    \multirow{2}*{\centering Phi-3 14B} 
       & 0\%      & 14B &  28 GB & 7171 t/s & 7216 t/s & 7197 t/s & 7057 t/s & 7010 t/s & 7036 t/s \\
       & 20\%     & 11B &  22.8 GB & 8622 t/s & 8686 t/s & 8509 t/s & 8349 t/s & 8268 t/s & 8269 t/s \\
    \midrule
       \multirow{2}*{\centering Mixtral-8x7B-v0.1 47B} 
       & 0\%      & 47B & OOM & OOM & OOM & OOM & OOM & OOM & OOM \\
       & 20\%     & 37B &  73.8 GB & 6515 t/s & 8143 t/s & 8444 t/s & OOM & OOM & OOM \\
    \midrule
       \multirow{2}*{\centering Phi-2 3B} 
       & 0\%      & 2.8B  &  6.8 GB & 29949 t/s & 30895 t/s & 30184 t/s & 30403 t/s & 26636 t/s & 21499 t/s \\
       & 20\%     & 2.2B  &  5.7 GB & 32451 t/s & 34561 t/s & 34276 t/s & 32479 t/s & 30351 t/s & 23735 t/s \\
    \midrule
       \multirow{2}*{\centering Mistral-v0.1 7B} 
       & 0\%      & 7.2B  &  15.2 GB & 13385 t/s & 13537 t/s & 13650 t/s & 13265 t/s & 12603 t/s & 12399 t/s \\
       & 20\%     & 5.8B  &  12.6 GB & 15416 t/s & 15832 t/s & 15947 t/s & 15446 t/s & 14589 t/s & 14304 t/s \\
    \bottomrule
    \end{tabular}}
    \caption{\textbf{Compressed models save memory and speed up computation.} Inference speed was measured on a single A100 GPU, using the whole test split of Wikitext2 with a batch size of 4 and varying the sequence length from 512 to 16384 tokens. All the models were loaded in bf16 and used Flash-Attention2. OOM stands for out of memory.}
    \label{tab:mem-speed}
\end{table*}

\begin{table*}[t]
    \resizebox{1.06\textwidth}{!}{%
    \setlength{\tabcolsep}{3pt}
    \begin{tabular}{lcccccccccccc}
        \toprule
        \textbf{Model} & \textbf{TruQA} & \textbf{SIQA} & \textbf{LogiQA} & \textbf{WinoG} & \textbf{ARC-E} & \textbf{ARC-C} & \textbf{BoolQ} & \textbf{PIQA} & \textbf{OBQA} & \textbf{Average} \\
        \midrule 
        StableLM-2 1.6B & 38.90 & 48.52 & 26.73 & 63.30 & 68.39 & 38.57 & 74.80 & 76.93 & 39.00 & \textbf{52.79} \\
        Qwen-1 1.8B & 38.09 &  45.19 & 31.80 & 59.12 & 58.33 & 34.98 & 65.93 & 73.23 & 33.60 & \textbf{48.92} \\
        Qwen-2 1.5B &  45.93 & 45.85 & 31.18 & 66.22 & 60.56 & 36.09 & 72.26 & 75.35 & 36.40 & \textbf{52.20} \\
        \textbf{Lillama} (Phi-2 1.7B) & 44.08 & 47.80 & 25.65 & 68.27 & 63.22 & 38.82 & 76.02 & 72.36 & 39.20 & \textbf{52.82} \\
        \bottomrule
    \end{tabular}}
    \caption{\textbf{40\% compressed Phi-2 3B competes with models of similar size while being compressed using only 13M tokens.} We compress Phi-2 3B to 1.7B (by removing 40\% of the parameters) and compare its zero-shot performance to that of other recent language models of similar sizes.}
    \label{tab:smollm}
\end{table*}

\begin{table*}[t]
    \resizebox{1.06\textwidth}{!}{%
    \setlength{\tabcolsep}{3pt}
    \begin{tabular}{cccccccccccc}
        \toprule
        \textbf{Model} & \textbf{Reduction} & \textbf{TruQA} & \textbf{SIQA} & \textbf{LogiQA} & \textbf{WinoG} & \textbf{ARC-E} & \textbf{ARC-C} & \textbf{BoolQ} & \textbf{PIQA} & \textbf{OBQA} & \textbf{Average} \\
        \midrule 
        \multirow{3}*{\centering Phi-2 3B} 
        & 0\% & 44.40 & 55.42 & 30.57 & 76.16 & 78.16 & 54.18 & 83.21 & 79.11 & 51.20 & \textbf{61.38} \\
        & 40\% & 44.08 & 47.80 & 25.65 & 68.27 & 63.22 & 38.82 & 76.02 & 72.36 & 39.20 & \textbf{52.82} \\
        & 40\% {\scriptsize FT} & 45.43 & 49.49 & 31.80 & 70.09 & 60.86 & 37.29 & 67.68 & 73.23 & 42.00 & \textbf{53.10} \\
        \midrule
        \multirow{3}*{\centering Mistral-v0.1 7B} 
        & 0\% & 42.60 & 46.57 & 29.80 & 73.80 & 79.55 & 53.92 & 83.70 & 82.10 & 44.00 & \textbf{59.56} \\
        & 40\% & 41.69 & 43.65 & 30.57 & 62.35 & 61.66 & 34.56 & 75.23 & 71.49 & 34.80 & \textbf{50.63} \\
        & 40\% {\scriptsize FT} & 41.60 & 52.25 & 29.65 & 66.61 & 66.75 & 39.93 & 79.39 & 74.37 & 38.29 & \textbf{54.32}\\
        \bottomrule
    \end{tabular}}
    \caption{\textbf{Compressed models recover well when fine-tuned.} Zero-shot performance of the base non-compressed models (0\%) and the 40\% compressed models, without and with fine-tuning (FT). Models were fine-tuned on the whole Slim-Orca dataset (191 million tokens).}
    \label{tab:ft_result}
\end{table*}


\section{Other Architectures and Modalities}
In this section, we show the generalizability of our approach by evaluating it on the Mamba architecture for text \citep{gu2024mambalineartimesequencemodeling} and on the speech modality with Whisper \citep{radford2022robustspeechrecognitionlargescale}.
We further evaluate it on a lower-resource African language later in this chapter.
\subsection{Mamba Architecture}

The attention mechanism in the Transformer is limited by its increasing complexity as a function of the input sequence length.
Therefore, Linear Attention Models have been proposed as alternatives to Transformers, in particular State Space Models, such as the Mamba architecture \citep{gu2024mambalineartimesequencemodeling}, which have shown promising results.

\paragraph{Method.} We tested our approach on two Mamba architectures: Mamba
3B \citep{gu2024mambalineartimesequencemodeling} (\url{https://huggingface.co/state-spaces/mamba-2.8b-hf}) and Falcon-Mamba 7B (\url{https://huggingface.co/tiiuae/falcon-mamba-7b}).
We compressed these models by 20\% using the same dataset and hyperparameters as in previous experiments with Transformers, with a minimum rank set to 1024 for both models.

\paragraph{Results.} Table~\ref{tab:ssm} shows that, as for Transformers, Mamba-based language models can be efficiently compressed. Interestingly, the 20\% compressed Mamba 3B maintains 99\% of the performance.
Table~\ref{tab:mem-speed-mamba-ch2} further shows that compressed Mamba models are lighter and faster than their original versions.
\begin{table*}[h!]
    \centering
    \resizebox{1.02\textwidth}{!}{%
    \setlength{\tabcolsep}{3pt}
    \begin{tabular}{lccccccccccc}
        \toprule
        \textbf{Model} & \textbf{Reduction} & \textbf{TruQA} & \textbf{SIQA} & \textbf{LogiQA} & \textbf{WinoG} & \textbf{ARC-E} & \textbf{ARC-C} & \textbf{BoolQ} & \textbf{PIQA} & \textbf{OBQA} & \textbf{Average} \\
        \midrule
        \multirow{2}*{\centering Falcon-Mamba 7B} 
        & 0\% & 53.40 & 52.66 & 31.34 & 74.66 & 81.73 & 58.96 & 83.12 & 81.56 & 48.60 & \textbf{62.89} \\
        & 20\% & 52.34 & 50.10 & 30.72 & 66.69 & 76.73 & 49.23 & 77.22 & 78.02 & 44.80 & \textbf{58.43} \\
        \midrule
        \multirow{2}*{\centering Mamba 3B} 
        & 0\% & 35.88 & 43.24 & 26.88 & 63.38 & 64.02 & 36.26 & 65.63 & 75.90 & 39.40 & \textbf{50.07} \\
        & 20\% & 37.50 & 42.43 & 27.04 & 61.40 & 62.16 & 35.58 & 64.89 & 73.29 & 40.00 & \textbf{49.37} \\
        \bottomrule
    \end{tabular}
    }
    \caption{\textbf{Mamba-based language models can be efficiently compressed}. Zero-shot evaluation results for the non-compressed (0\%) and for the compressed Mamba-based language models (20\%).}
    \label{tab:ssm}
\end{table*}

\begin{table*}
    \centering
    \resizebox{\linewidth}{!}{%
    \setlength{\tabcolsep}{3pt}
    \begin{tabular}{ccccc|ccccc}
    \toprule
       \textbf{Model} & \textbf{Reduction} & \textbf{Model Size} & \textbf{VRAM} & s = 512 & s = 1024 & s = 2048 & s = 4096 & s = 8192 & s = 16384 \\
    \midrule
    \multirow{2}*{\centering Falcon-Mamba 7B}
       & 0\%      & 7.27B & 15.3 GB & 8725 t/s & 9346 t/s & 9390 t/s & 9695 t/s & 9908 t/s & 9486 t/s \\
       & 20\%     & 5.82B & 12.6 GB & 10227 t/s & 10427 t/s & 10477 t/s & 10814 t/s & 11108 t/s & 10647 t/s \\
    \midrule
    \multirow{2}*{\centering Mamba 3B}
       & 0\%      & 2.77B & 6.7 GB & 17994 t/s & 18556 t/s & 18819 t/s & 19618 t/s & 19231 t/s & 18602 t/s \\
       & 20\%     & 2.22B & 5.6 GB & 19013 t/s & 19993 t/s & 20378 t/s & 21218 t/s & 20808 t/s & 20100 t/s \\
    \bottomrule
    \end{tabular}}
    \caption{\textbf{Compressing Mamba models saves memory and speeds up computation.} Inference speed was measured on a single A100 GPU, using the whole test split of Wikitext2 with a batch size of 4 and by varying the sequence length from 512 to 16384 tokens. All models were loaded in bf16 and used Flash-Attention2.}
    \label{tab:mem-speed-mamba-ch2}
\end{table*}


\subsection{Speech modality: Whisper}

We next evaluate our approach on a state-of-the-art speech model: \textit{whisper-medium.en} \citep{radford2022robustspeechrecognitionlargescale}, a 764M-parameter encoder-decoder speech recognition model trained on English.

\paragraph{Method.} We take inspiration from the deep encoder,
shallow decoder approach \citep{kasai2021deepencodershallowdecoder}, where the
authors observe that the decoder in encoder-decoder models does not need to be very deep to achieve good performance.
Therefore, we keep the encoder unchanged and remove 70\% of the decoder's parameters. We then approximate the low-rank matrices using the same method and hyperparameters described in Section \ref{sec:xp}.
\paragraph{Data.} For calibration, we used a random 60\% sample of examples from Librispeech-train-100 \citep{librispeech}. We also fine-tuned the compressed model on a subset of 10 hours of speech randomly sampled from Librispeech-train-100.

\paragraph{Results.} We evaluate the models' Word Error Rate (WER), memory footprint and latency on an out-of-distribution test dataset: the test split of the English FLEURS dataset \citep{fleurs}.

\begin{table}[!ht]
    {\small{
    \begin{tabular}{ccccc}
    \toprule
        \textbf{Reduction} &\textbf{Size} & \textbf{WER ($\downarrow$)}  & \textbf{VRAM} & \textbf{Speed} ($\uparrow$) \\
    \midrule
        0\% & 764M & 25.40 & 3.5GB & 2395t/s\\
        37\%  & 485M & 30.84 & 2.4GB & 2717t/s \\
        0\% {\scriptsize FT} & 764M & 9.20 & 3.5GB & 2395t/s \\
        37\% {\scriptsize FT} & 485M & 12.14 & 2.4GB & 2717t/s \\ 
    \bottomrule
    \end{tabular}
    }}
    \caption{\textbf{The 37\% compressed Whisper is faster and 1.46$\times$ lighter.} The WER is computed on the English test subset of the FLEURS dataset. The speed (tokens/second) corresponds to the transcription of 400 tokens from a dataset containing 2,048 speech examples, using a batch size of 128. A single A100 GPU is used and the model is loaded in FP32, without Flash-Attention. We also give the GPU Memory (VRAM) occupied by each model.}
    \label{tab:wer}
\end{table}

Table~\ref{tab:wer} gives the WER for compressed and uncompressed models. We can see that a 37\% compressed Whisper maintains 82\% of the performance.

\section{Analysis}

\begin{figure}[ht]
    \centering
    \includegraphics[width=\linewidth]{papers/lillama/rank_acc-1.pdf}
    \caption{\textbf{Higher ranks generally give better results.} Mean accuracy over our test benchmarks as a function of the minimum rank $k$ in Algorithm~\ref{alg:two}}
    \label{fig:ranks}
\end{figure}

\begin{figure}[ht]
    \centering
    \includegraphics[width=\linewidth]{papers/lillama/w_wo_svd-2.pdf}
    \caption{\textbf{SVD initialization improves sample efficiency}. Convergence when initializing low-rank weights randomly (without SVD) or with SVD.}
    \label{fig:w_wo_svd}
\end{figure}


In this section, we give some ablation studies. In all experiments, the models are compressed by 20\%.


\subsection{Which Distillation Loss Should Be Used?}
\begin{figure*}[h!]
    \centering
        \includegraphics[width=1.0\textwidth]{papers/lillama/losses_convergence.pdf}
        \caption{\textbf{The joint loss generally converges better}. Convergence of the three losses illustrated in Figure~\ref{fig:distillaton}, evaluated using perplexity on the Wikitext2 test corpus during distillation.}
        \label{fig:losses}
\end{figure*}

\begin{table}[h!]
    \centering
    \begin{tabular}{lccc}
    \toprule
    \textbf{Model} & \textbf{Teacher} & \textbf{Student} & \textbf{Tea+Stu} \\
    \midrule
    Phi-2 3B & 57.38 & 56.44 & 57.38 \\
    Phi-3 14B & 62.54 & 63.20 & 63.36 \\
    Mistral-v0.1 7B & 51.25 & 55.46 & 55.51 \\
    Mixtral-8x7B 47B & 58.87 & 60.49 & 60.19 \\
    \midrule
    \textbf{Average} & 57.51 & 58.90 & \textbf{59.11} \\
    \bottomrule
    \end{tabular}
    \caption{\textbf{The joint loss generally performs better}: Zero-shot performance of the three losses using the same setup for each model, as described in Section~\ref{sec:xp}.}
    \label{tab:losses}
\end{table}

Table~\ref{tab:losses} shows that the joint loss depicted in Figure~\ref{fig:distillation_joint} generally produces the best results. Although the differences between the Student and Teacher+Student losses may not seem significant, Figure~\ref{fig:losses} shows that the Teacher+Student loss leads to faster convergence than the Student loss alone. It also generally results in a lower final perplexity than using the Teacher loss only.

\begin{insightbox}{Observation 3: Combining Teacher and Student activations improves average performance}
Combining Teacher and Student activations in a joint loss achieves the highest average zero-shot score across the evaluated models (Table~\ref{tab:losses}). It also converges faster than the Student-only loss and generally reaches a lower final perplexity than the Teacher-only loss (Figure~\ref{fig:losses}).
\end{insightbox}

As shown in Figure~\ref{fig:losses}, the Teacher+Student loss combines the best of both worlds: fast convergence due to the Teacher activations and high performance at inference thanks to the Student activations. We illustrate this in Table~\ref{tab:logs} for the first 2M training tokens. It shows that the Teacher+Student loss converges faster than the Student-only loss. While the Teacher-only loss also converges fast, it reaches a higher plateau and does not converge beyond that point.

\begin{table}[!ht]
    {\scriptsize{
    \begin{tabular}{lcccc}
    \toprule
    \textbf{Model} & \textit{Seen Tokens} & \textbf{Teacher} & \textbf{Student} & \textbf{Tea+Stu} \\
    \midrule
    \multirow{5}{*}{Phi-2 3B}
                              & 0       & 7185.31  & 7185.31  & 7185.31 \\
                              & 532480  & 20.74    & 36.56    & 20.34 \\
                              & 1056768 & 17.69    & 25.11    & 17.26 \\
                              & 1559879 & 16.00    & 18.97    & 15.92 \\
                              & 2000506 & 15.21    & 18.19    & 15.51 \\
    \midrule
    \multirow{5}{*}{Mistral-v0.1 7B}
                              & 0       & 12365.44 & 12365.44 & 12365.44 \\
                              & 532480  & 43.79    & 270.16   & 114.02 \\
                              & 1056768 & 19.07    & 53.90    & 21.90 \\
                              & 1579121 & 18.56    & 22.40    & 13.64 \\
                              & 2056609 & 18.55    & 15.72    & 13.11 \\
    \midrule
    \multirow{5}{*}{Phi-3 14B}
                              & 0       & 46362.82 & 46362.82 & 46362.82 \\
                              & 532480  & 9.77     & 12.94    & 8.89 \\
                              & 1056768 & 8.33     & 8.76     & 7.99 \\
                              & 1580799 & 7.70     & 7.85     & 7.44 \\
                              & 2069488 & 7.61     & 7.77     & 7.34 \\
    \bottomrule
    \end{tabular}
    }}
    \caption{Comparing the Wikitext perplexity convergence of the three losses for 2M training tokens.}
    \label{tab:logs}
\end{table}

\subsection{How Should the Rank Be Chosen?}
In our bottom-first compression strategy, a small minimum rank in
Algorithm~\ref{alg:two} compresses the bottom layers more aggressively and leaves
more upper layers uncompressed, hence saving GPU memory. However, a
small rank may also destroy knowledge and the model will struggle to recover
during the distillation phase. Figure~\ref{fig:ranks} compares 4 rank
values for Mistral-v0.1 7B and Phi-3 14B and shows that higher ranks
lead to better performance, but are also more costly as more layers are distilled.

\subsection{Is SVD initialization necessary?}
Figure~\ref{fig:w_wo_svd} shows the perplexity curves for Phi-2 3B and Phi-3 14B when the low-rank matrices are
initialized with SVD or randomly. SVD initialization enables faster
convergence, with 8 million tokens being sufficient, i.e., roughly
4k examples of 2048 tokens.

\subsection{Cost.}
All models in Section~\ref{sec:xp} are compressed in less than one hour on a single A100 GPU. The precise cost depends on the number of layers distilled, which is determined by the compression strategy (bottom-first, top-first, or uniform) and the minimum rank $k$ in Algorithm~\ref{alg:two}. Figure~\ref{fig:strategies} intuitively illustrates the three compression strategies evaluated in our method: uniform, bottom-first, and top-first, explaining their respective advantages and drawbacks. We present here empirical results on differences in efficiency and performance using the same hyperparameters as in Section~\ref{sec:xp}.

\paragraph{Differences in running time and memory.}
The bottom-first compression strategy allows compressing all models, including Mixtral-8x7B-v0.1 47B. While the top-first strategy can compress Phi-3 14B, the uniform compression strategy fails to compress Phi-3 14B and Mixtral-8x7B-v0.1 (47B). Table~\ref{tab:ctime-ch2} gives the distillation time for each strategy and each model.

\begin{table}[!ht]
    {\small{
    \begin{tabular}{lccc}
    \toprule
    \textbf{Model} & \textbf{Bottom} & \textbf{Top} & \textbf{Uniform} \\
    \midrule
    Phi-2 3B & 25min & 32min & 50min \\
    Mistral-v0.1 7B & 26min & 46min & 105min \\
    Phi-3 14B & 30min & 50min & OOM \\
    Mixtral-8x7B-v0.1 47B & 47min & OOM & OOM \\
    \bottomrule
    \end{tabular}
    }}
    \caption{Compression times for the three compression strategies at 20\% parameter reduction on a single A100 GPU. For the bottom-first and top-first compression strategies, we used the same hyperparameters as in Section~\ref{sec:xp}.}
    \label{tab:ctime-ch2}
\end{table}

\begin{insightbox}{Observation 4: Bottom-first compression is faster and more memory efficient}
Bottom-first compression reduces computation and GPU memory requirements by processing only the required bottom layers. It is the fastest strategy across all evaluated models (Table~\ref{tab:ctime-ch2}). It is also the only strategy that compresses Mixtral 47B on a single A100 GPU, taking 47 minutes.
\end{insightbox}

\paragraph{Differences in performance.}
Table~\ref{tab:cacc-ch2} presents the mean accuracies for each compression strategy. The bottom-first compression strategy performs well and is also more memory efficient. The top-first strategy does not improve accuracy while having higher time and memory complexity compared to the bottom-first approach.

\begin{table}[!ht]
    {\small{
    \begin{tabular}{lccc}
    \toprule
    \textbf{Model} & \textbf{Bottom} & \textbf{Top} & \textbf{Uniform} \\
    \midrule
    Phi-2 3B & 57.38 & 55.71 & 58.21 \\
    Mistral-v0.1 7B & 55.51 & 55.63 & 54.25 \\
    Phi-3 14B & 63.59 & 60.80 & OOM \\
    Mixtral-8x7B-v0.1 47B & 60.19 & OOM & OOM \\
    \bottomrule
    \end{tabular}
    }}
    \caption{Average zero-shot performance for the three compression strategies at 20\% compression. For the bottom-first and top-first compression strategies, we used the same hyperparameters as in Section~\ref{sec:xp}.}
    \label{tab:cacc-ch2}
\end{table}

\begin{table*}[h!]
    \centering
    \resizebox{\textwidth}{!}{%
    \begin{tabular}{cccccccc}
        \toprule
        \textbf{Reduction} & \textbf{Method} & \textbf{PIQA} & \textbf{WinoG} & \textbf{HellaSwag} & \textbf{ARC-E} & \textbf{ARC-C} & \textbf{Avg.} \\
        \midrule
        \multirow{2}{*}{0\%}
        & SliceGPT & 79.11 & 75.77 & 73.83 & 78.32 & 54.18 & 72.24 \\
        & Reproduced & 79.11 & 76.01 & 73.60 & 78.41 & 54.35 & 72.30 \\
        \midrule
        \multirow{2}{*}{24\%}
        & SliceGPT & 74.05 & 62.12 & 53.31 & 67.26 & 39.42 & 59.23 \\
        & \textbf{Lillama} & \textbf{74.76} & \textbf{69.77} & \textbf{60.60} & \textbf{72.10} & \textbf{47.10} & \textbf{64.86} \\
        \bottomrule
    \end{tabular}
    }
    \caption{\textbf{Comparison of our method with SliceGPT on zero-shot task evaluation for Phi-2 3B.} We give the evaluation scores before compression (0\%), both as reported by SliceGPT and as reproduced by us. Both methods use the Alpaca calibration dataset for compression.}
    \label{tab:comparisons}
\end{table*}

\subsection{Examples of generated texts}
We also evaluated the compressed models by observing the text they generated. In general, at 20\% compression of Phi-2 3B, the generation ability and knowledge remained nearly intact. At 40\% compression, the model still retained its general knowledge and generation abilities but tended to be more repetitive and verbose at times. We recommend users fine-tune the 40\% compressed Phi-2 3B on their specific tasks, such as retrieval-augmented generation. Tables~\ref{tab:sj} and~\ref{tab:cw} give some generated texts from 20\% and 40\% compressed Phi-2 3B.

\begin{table*}
\scriptsize
\centering
\setlength{\tabcolsep}{3pt}
\begin{tabular}{ p{0.31\textwidth} | p{0.31\textwidth} | p{0.31\textwidth} }
\toprule
Base Model & -20\% & -40\% \\
\midrule
    Steve Jobs was a young man who had a passion for technology and design. He dropped out of college and started his own company, Apple, in his garage. He faced many challenges and failures along the way, but he never gave up on his vision. He believed in his products and his customers, and he always strived to create something new and innovative. He was not afraid to take risks and to challenge the status quo. He was a leader who inspired others to follow his example. He was a visionary who changed the world. &  Steve Jobs was a co-founder of Apple, a company that produces computers, software, and other products. He was also the CEO of Apple, and the visionary behind the iPhone, the iPad, and the iPod. He was known for his innovation, creativity, and leadership. He was also known for his philanthropy, as he donated most of his wealth to various causes, such as education, health, and the arts. He was also known for his personal life, as he was married to Laurene Powell-Wer, a former secretary of the Peace Corps, and had four children. He died in 2011 at the age of 56, after suffering from pancreatic cancer.

Questions:

- What was Steve Jobs' occupation?
- What were some of the products that Apple produced?
- What were some of the causes that Steve Jobs donated to?
- How many children did Steve Jobs have?
- When and how did Steve Jobs die?

Answers:

- Steve Jobs' occupation was a co-founder of Apple and the CEO of Apple.
- Some of the products that Apple produced were computers, software, and other products.
- Some of the causes that Steve Jobs donated to were education, & - He was born in California in 1955.

- He dropped out of college and started working at a computer company.

- He became the co-founder of Apple.

- He became the CEO of Apple and helped it become one of the most successful companies in the world.

- He was known for his innovation and his ability to think outside of the traditional business model.

- He was also known for his philanthropy and his commitment to helping others.

- He died in 2011.
\\
\bottomrule
\end{tabular}
\caption{Steve Jobs Story by the 20\% and 40\% compressed Phi-2 3B models, without \textit{any} fine-tuning.
\textbf{Prompt: "Here the story of Steve Jobs:"}}
\label{tab:sj}
\end{table*}

\begin{table*}
\scriptsize
\centering
\setlength{\tabcolsep}{3pt}
\begin{tabular}{ p{0.31\textwidth} | p{0.31\textwidth} | p{0.31\textwidth} }
\toprule
Base Model & -20\% & -40\% \\
\midrule
The Civil War was caused by a number of factors, including economic and social differences between the North and the South, as well as disagreements over states' rights and slavery. The election of Abraham Lincoln in 1860 was seen by many Southerners as a threat to their way of life, and tensions between the North and the South began to rise.
The war began in 1861 when Confederate forces attacked Fort Sumter in South Carolina. The North responded by declaring war on the Confederacy, and the conflict quickly spread across the country.
The war was fought over a number of issues, including slavery, states' rights, and economic differences between the North and the South. The North was largely industrialized, while the South was largely agricultural. This led to tensions over tariffs and other economic issues.
The war was also fought over the issue of slavery. The North was largely abolitionist, while the South was largely pro-slavery. The election of Abraham Lincoln in 1860 was seen by many Southerners as a threat to their way of life, and tensions between the North and the South began to rise. & The American Civil War was a conflict that took place from 1861 to 1865. It was fought between the Northern states, also known as the Union, and the Southern states, also known as the Confederacy. The main cause of the war was the issue of slavery. The Southern states wanted to keep slavery, while the Northern states wanted to abolish it. This led to tensions and eventually led to the war.

What were the major battles of the Civil War? Here is the story.

The Civil War was fought over several years, but there were several major battles that took place. The most famous battles were the Battle of Gettysburg, the Battle of Antieta, and the Battle of Bull Run. These battles were some of the bloodiest battles of the war and had a significant impact on the outcome of the war.

What was the Emancipation Proclamation? Here is the story.

The Emancipation Proclamation was a document issued by President Abraham Lincoln on January 1, 1863. It declared that all slaves in Confederate territory were to be set free. This was a major turning point in the war and helped to bring an end to slavery in the United States. & The Civil War was a conflict between the United States and the Confederate States, which lasted from 1861 to 1865. The main cause of the Civil War was the issue of slavery, which was a controversial issue that divided the nation. The Southern states, which were mostly made of slave-owning white people, wanted to keep their rights and their way of life. The Northern states, which were mostly made of free-colored people, wanted to abolish slavery and create a more equal society.

The Civil War began when the Confederate States declared their secession from the United States, which means they wanted to break away from the country. The United States responded by sending troops to the South, which started the war. The war lasted for four years, and it was one of the most important events in American history. The war ended with the defeat of the Confederate States, and the United States became a unified nation.

The Civil War had many consequences, both positive and negative. Some positive consequences were that the United States became more united and more democratic, and that the issue of slavery was finally resolved. Some negative consequences were that many people died, and that the war caused a lot of damage and destruction.
\\
\bottomrule
\end{tabular}
\caption{The Civil War Story by the 20\% and 40\% compressed models, without \textit{any} fine-tuning. \textbf{Prompt: "What is the cause of the Civil War? Here is the story."}}
\label{tab:cw}
\end{table*}

\subsection{Comparisons.}
Table~\ref{tab:comparisons} shows that our approach outperforms SliceGPT \citep{ashkboos_slicegpt_2024}, a one-shot compression method for LLMs.
SliceGPT reports a ``slicing'' proportion, which does not directly correspond to the percentage reduction in the model's parameter count.
For example, a 30\% slicing proportion for Phi-2 3B corresponds to a 24\% reduction in its parameter count.

\subsection{Can this approach work for low-resource and under-represented languages?}
\label{sec:hausa-ch2}

For most human languages, there is not enough data (generally less than 1B tokens) to pretrain large LLMs. In such cases, data-efficient compression approaches are preferred. We experimented with compressing a tiny language model for Hausa, a low-resource language for which there are not billions of tokens available for continued pretraining. We used InkubaLM~\citep{tonja2024inkubalm}, a 422M-parameter language model designed for five low-resource languages. Since 60\% of the parameters are in the embeddings, we focused on compressing the input embedding $W_{e}$ and the prediction head $W_{o}$. We will discuss later on the use of low-rank structure in the LM head, as it appears to contradict some empirical observations from previous studies. For now, let's focus on compressing the input and output modules by applying our approach through local distillation of the low-rank embeddings and prediction head:

\[
\widehat{\Delta W} = \underset{{\Delta W}}{\mathrm{argmin}} \;\; \mathcal{L}(Wx, \Delta Wx)
\]
where $\Delta W$ are the low-rank embedding matrices or prediction head, and $x$ is an input example.

We used a rank of 1024, compressing the model by 30\%. Then, we applied lightweight local distillation using our approach, with 64k randomly sampled Hausa sentences from InkubaMono\footnote{\url{https://huggingface.co/datasets/lelapa/Inkuba-Mono}} and the Aya-Dataset~\citep{singh2024aya}. We compared the models with MobiLLaMA~\citep{thawakar2024mobillama} and SmoLLM\footnote{\url{https://huggingface.co/blog/smollm}} using the AfriMMLU benchmark~\citep{afrimmlu}. As shown in Table~\ref{tab:inkuba-ch2}, the compressed InkubaLM model for Hausa retains 93\% of its base performance.

\begin{table}[h!]
    \centering
    {\small{
    \begin{tabular}{lcc}
    \toprule
    \textbf{Model} & \textbf{Model Size (Billion)} & \textbf{Accuracy} \\
    \midrule
    SmoLLM & 1.7 & 21.80 \\
    MobiLLaMA & 1.3 & 21.40 \\
    \midrule
    InkubaLM-base & 0.422 & 29.20 \\
    InkubaLM-30\% & 0.299 & 27.40 \\
    \bottomrule
    \end{tabular}
    }}
    \caption{Zero-shot performance on AfriMMLU for the 30\% compressed InkubaLM model for the Hausa language.}
    \label{tab:inkuba-ch2}
\end{table}

\section{Limitations}

\paragraph{Compatibility with quantization:}
Our pruning/compression method is theoretically complementary to quantization, since a model compressed with our approach can be further quantized. However, each step removes information, making it crucial to balance these approaches optimally. Addressing this challenge is essential for large-scale adoption in LLM libraries but is beyond the scope of this work and warrants a dedicated study.

\paragraph{Compromise between specificity and generality:}
Efficient LLM compression requires trade-offs between speed and generality. We prioritize generality over raw speed, avoiding GPU-specific kernels as used in semi-structured sparsity approaches. Efficiency also ties to carbon cost: our method works on CPUs as well as on
GPUs, enabling the reuse of older hardware and reducing reliance on GPUs with higher carbon costs.

\section{Conclusion}
We propose a new efficient compression method that
requires far less calibration data than most state-of-the-art approaches
and provides competitive performance for various models (dense and MoE Transformers, Mamba) and modalities (text, speech).
We show that combining SVD initialization with a joint Teacher and Student loss for local distillation enables fast convergence, and a bottom-up layer selection approach enables cost-effective compression.
